The fundamental domain describes the quotient through its boundary identifications. We now determine the stabilizers at the identified points and the geometry of its single end. Together, these show that the underlying space is a punctured sphere, while its orbifold structure records two elliptic points.

\begin{proposition}[Elliptic stabilizers]\label{modular.elliptic-stabilizers}
In $\mathcal F$, the only points with nontrivial effective stabilizer are $i,\rho,\rho+1$. The points $\rho$ and $\rho+1$ are equivalent, and
\[
 \begin{aligned}
 \Gamma_i&=\langle S\rangle,& |\Gamma_i|&=4,& |\overline\Gamma_i|&=2,\\
 \Gamma_\rho&=\langle ST\rangle,& |\Gamma_\rho|&=6,& |\overline\Gamma_\rho|&=3.
 \end{aligned}
\]
At every other point not equivalent to these, the matrix stabilizer is $Z$ and the effective stabilizer is trivial. Thus the quotient has exactly two elliptic point orbits.
\end{proposition}

\begin{proof}
An interior stabilizer is central by the \hyperref[modular.fundamental-domain]{fundamental-domain theorem}. On a vertical side away from the endpoints, the only possible noncentral equivalence is translation to the other side, which does not fix the point. On the circular arc away from the endpoints, the only possible nonidentity effective stabilizer is $S$, and $Sz=z$ there exactly when $z=i$. Directly, a matrix fixing $i$ has $a=d,b=-c$, and $a^2+c^2=1$ with integer entries. The four solutions are $\pm I,\pm S$.

At $\rho$, the equation $\gamma\rho=\rho$ and $\rho^2+\rho+1=0$ give $b=-c$, $d=a+c$. The determinant becomes $a^2+ac+c^2=1$. Since $(a+c/2)^2+3c^2/4=1$, one has $c=0,1,-1$; these give exactly six integer solutions, the powers of $ST$. Translation conjugates this stabilizer to the one at $\rho+1$. The relation $(ST)^3=-I$ gives effective order three. Dividing each matrix stabilizer by $Z$ gives the stated orders and completes the classification.
\end{proof}

\begin{remark}[Local quotient geometry]\label{modular.orbifold}
Near an elliptic point with effective stabilizer of order $m$, move the point to the center of the disk. Its stabilizer is a finite subgroup of rotations, hence is generated by a rotation through $2\pi/m$: a finite subgroup of the circle consists of equally spaced roots of unity, as follows by taking its smallest positive angle. Properness permits a disk small enough that only stabilizer elements meet it with their translates. To see this, begin with a compact neighborhood, use finiteness of intersecting translates, and shrink to avoid each of the finitely many elements not fixing its center. The local quotient identifies angles modulo $2\pi/m$ and has cone angle $2\pi/m$. An invariant integral over this quotient neighborhood is $1/m$ times the integral of its lift over the whole small disk. This is the local stabilizer factor, not an additional global rescaling of the fundamental-domain area. The modular quotient is an orbifold with these two cone points.
\end{remark}

\begin{definition}[Cusps and the stabilizer at infinity]\label{modular.cusps}
A cusp of a discrete subgroup is an orbit of boundary points fixed by a nonidentity parabolic transformation. For the modular group, write $\Gamma_\infty$ for the full matrix stabilizer of $\infty$ and $\overline\Gamma_\infty$ for its effective image.
\end{definition}

\phantomsection\label{modular.horocycle}
For $Y>0$, the horizontal line $y=Y$ is called a horocycle based at infinity. The horizontal-segment length formula \eqref{eq:horizontal-segment-length} gives hyperbolic length $1/Y$ to a segment of Euclidean width one. The \hyperref[modular.cusp-coordinate]{cusp-coordinate proposition} explains why identifying $z$ with $z+1$ makes such a line a closed cross-section of the modular cusp neighborhood.

\begin{proposition}[The single modular cusp]\label{modular.cusp-stabilizer}
The parabolic fixed points of $\Gamma$ are exactly $\mathbb Q\cup\{\infty\}$, and they form one orbit. Moreover,
\[
 \Gamma_\infty=\{\pm T^n:n\in\mathbb Z\},\qquad
 \overline\Gamma_\infty=\langle\overline T\rangle.
\]
\end{proposition}

\begin{proof}
An integer matrix fixes infinity precisely when $c=0$. The determinant then forces $a=d=\pm1$, giving the stated stabilizer. For a noncentral parabolic integer matrix, the fixed-point quadratic has discriminant zero. If $c\ne0$ its root is $(a-d)/(2c)\in\mathbb Q$; if $c=0$ its unique boundary fixed point is infinity. Conversely, a reduced rational $a/c$ is the image of infinity under a determinant-one integer matrix: choose $b,d$ with $ad-bc=1$. Conjugating $T$ by that matrix supplies a parabolic fixing $a/c$, proving both assertions about cusps.
\end{proof}

\begin{definition}[Cusp width in a chosen coordinate]\label{modular.cusp-width}
Suppose a cusp has been moved to $\infty$ by a real fractional linear transformation, and in this coordinate its effective stabilizer consists of the translations
\[
 z\longmapsto z+nh,\qquad n\in\mathbb Z,
\]
for some $h>0$. The number $h$ is the width of the cusp in this coordinate: it is the least positive horizontal translation in the stabilizer. Thus a horizontal interval of Euclidean width $h$ is one full period of the identification.
\end{definition}

\begin{remark}[The width of the modular cusp]\label{modular.cusp-width-one}
By \hyperref[modular.cusp-stabilizer]{Proposition 1.7.4}, the effective stabilizer of $\infty$ for $\Gamma=\operatorname{SL}_2(\mathbb Z)$ consists of the translations $z\mapsto z+n$, $n\in\mathbb Z$. Its least positive translation is $1$, so the cusp at infinity for $\Gamma$ has width $1$ in the original coordinate $z$.
\end{remark}

The coordinate normalization matters: rescaling $z$ by a positive constant rescales $h$ by the same constant. Cusp width is a translation period; the horizontal-segment length formula \eqref{eq:horizontal-segment-length} gives hyperbolic circumference $h/Y$ at height $Y$.

\begin{definition}[Exponential notation and the cusp parameter]\label{modular.cusp-parameter}
For $z\in\mathbb C$, write
\[
 e(z)=\exp(2\pi iz).
\]
On $\mathbb H$ define the cusp parameter $q(z)=e(z)$. When the point $z$ is understood, we write $q=q(z)$.
\end{definition}

The parameter $q$ records a complex number associated to the point $z$; it is not a second independent variable. If $z=x+iy$, the exponential formula gives
\begin{equation}\label{eq:cusp-parameter-polar}
 q(x+iy)=e^{-2\pi y}\bigl(\cos(2\pi x)+i\sin(2\pi x)\bigr),
 \qquad |q(x+iy)|=e^{-2\pi y}.
\end{equation}
Thus $x$ determines the angle modulo one full turn, while $y$ determines the radius. In particular, $q(z+1)=q(z)$: identifying the vertical sides of a width-one strip makes its horizontal lines into circles. As $y\to\infty$, these circles shrink toward $q=0$, which is never attained at a finite point of $\mathbb H$. For a translation period $h$ the corresponding parameter would be $e(z/h)$; here the modular width is $1$.

\begin{proposition}[A coordinate on the modular cusp neighborhood]\label{modular.cusp-coordinate}
Let $Y>1$. The image of $\{z\in\mathbb H:\operatorname{Im}z>Y\}$ in $\overline\Gamma\backslash\mathbb H$ is an open neighborhood of the cusp end. Two points of this region represent the same quotient point exactly when they differ by an integer. The map
\[
 [z]\longmapsto q(z)
\]
is a homeomorphism from this neighborhood to the punctured disk $\{q\in\mathbb C:0<|q|<e^{-2\pi Y}\}$, where $[z]$ denotes the quotient point represented by $z$.
\end{proposition}

\begin{proof}
First determine which points are identified. If $z=x+iy$ and $\gamma\cdot z$ both have height greater than $Y$ and the bottom-left entry $c$ of $\gamma$ is nonzero, then the imaginary-part formula \eqref{eq:mobius-imaginary-part} gives
\[
 \operatorname{Im}(\gamma\cdot z)=\frac{y}{|cz+d|^2}
 \leq\frac1{c^2y}<\frac1Y<Y,
\]
a contradiction. Hence $c=0$, and the \hyperref[modular.cusp-stabilizer]{stabilizer calculation} gives $\gamma=\pm T^n$. Conversely, these matrices identify $z$ with $z+n$ and preserve its height. The image of the region is open in the quotient because its inverse image is the union of its open translates.

Next check that $q$ gives exactly one coordinate for each such quotient point. By the polar formula \eqref{eq:cusp-parameter-polar}, equality $q(z_1)=q(z_2)$ forces equal radii, hence equal heights. Equality of the angular factors then gives $\operatorname{Re}z_1-\operatorname{Re}z_2\in\mathbb Z$. The converse follows from period one. Therefore
\begin{equation}\label{eq:cusp-parameter-fibers}
 q(z_1)=q(z_2)\quad\Longleftrightarrow\quad z_1-z_2\in\mathbb Z.
\end{equation}
Here the right side views $\mathbb Z$ as a subset of $\mathbb C$. The fiber identity \eqref{eq:cusp-parameter-fibers} shows that $[z]\mapsto q(z)$ is well-defined and injective. The radius formula \eqref{eq:cusp-parameter-polar} puts its image in the stated punctured disk.

To prove surjectivity, take a point $q=re^{i\theta}$ with $0<r<e^{-2\pi Y}$ and $\theta\in\mathbb R$, and set
\begin{equation}\label{eq:cusp-coordinate-inverse}
 z=\frac{\theta}{2\pi}-\frac{i}{2\pi}\log r.
\end{equation}
Its height is $-\log r/(2\pi)>Y$, and substitution gives $q(z)=q$. Choosing another argument $\theta+2\pi n$ changes $z$ to $z+n$, which represents the same quotient point.

Finally, the exponential is continuous, so the descended map is continuous. Near any nonzero point of the disk, choose a continuous branch of the argument. The inverse-coordinate formula \eqref{eq:cusp-coordinate-inverse} then gives a continuous local inverse, and different branches agree in the quotient. Thus the bijection is a homeomorphism. The restriction of the modular quotient topology agrees with this translation quotient topology: for an open subset of the region, the full modular saturation is open, and its intersection with the region consists exactly of its integer translates.
\end{proof}

\begin{corollary}[Length and area of the modular cusp]\label{modular.cusp-area}
For $Y>1$, the closed horizontal cross-section at height $Y$ has hyperbolic length $1/Y$, and the cusp neighborhood above height $Y$ has hyperbolic area $1/Y$.
\end{corollary}

\begin{proof}
The \hyperref[modular.cusp-coordinate]{cusp-coordinate proposition} identifies the two ends of a horizontal segment of Euclidean width one, so the horizontal-segment length formula \eqref{eq:horizontal-segment-length} gives its length $1/Y$. For the cross-section at height $Y$, the same identification statement applies using any cutoff strictly between $1$ and $Y$. Integrating over a width-one strip gives
\[
 \int_Y^\infty\int_{-1/2}^{1/2}\frac{dx\,dy}{y^2}
 =\int_Y^\infty\frac{dy}{y^2}
 =\frac1Y.
\]
\end{proof}

\begin{theorem}[The topology of the modular quotient]\label{modular.quotient-topology}
The underlying topological space of $\overline\Gamma\backslash\mathbb H$ is a sphere with one point removed. The missing point is the cusp. As a hyperbolic orbifold, the quotient also has two cone points, of orders $2$ and $3$, represented by $i$ and $\rho$.
\end{theorem}

\begin{proof}
First add the cusp point. The \hyperref[modular.cusp-coordinate]{cusp-coordinate homeomorphism} identifies its punctured neighborhood with $0<|q|<e^{-2\pi Y}$; adjoining $q=0$ fills this puncture. The compactified fundamental domain is obtained by adjoining one point at the upper end of $\mathcal F$. To see its shape, write
\[
 y=\sqrt{1-x^2}+u,\qquad -\tfrac12\leq x\leq\tfrac12,\quad u\geq0.
\]
The parameter $u/(1+u)$ turns this into a rectangle with its top edge omitted. Adjoining the cusp collapses that top edge to a single vertex. Thus the compactified domain is a closed topological disk with vertices $\rho,\rho+1,\infty$.

Cut this disk along the imaginary axis from $i$ to the cusp, as indicated by the dashed line in the \hyperref[modular.domain-diagram]{domain diagram}. The two resulting closed disks have vertices $\rho,i,\infty$ and $\rho+1,i,\infty$. Their boundaries are paired in three pieces: $S$ identifies their circular arcs, $T$ identifies their vertical sides, and the identity rejoins the two copies of the cut. These pairings agree at the three vertices and together form a homeomorphism between the entire boundary circles. There are no other identifications, by the \hyperref[modular.boundary-identifications]{boundary-identification proposition}.

Two closed disks glued along their boundary circles by a homeomorphism form a sphere. Indeed, extend that boundary homeomorphism radially across one disk; the gluing then becomes the identity gluing of two disks, which can be realized as the northern and southern hemispheres. This proves that the compactified quotient is a sphere. The side-pairing description gives the quotient topology: away from the cusp this follows from properness and the local quotient neighborhoods, while at the cusp it follows from the coordinate $q$. Removing the single added point gives the asserted punctured sphere.

Finally, the \hyperref[modular.elliptic-stabilizers]{stabilizer calculation} and \hyperref[modular.orbifold]{local rotation model} give cone orders $2$ and $3$. These points are present in the quotient; they are not punctures. A cone neighborhood has underlying topology a disk (in a disk coordinate, the rotation quotient is represented by $w\mapsto w^m$), but retains the hyperbolic cone angle $2\pi/m$. Thus the sphere description and the orbifold description record different aspects of the same quotient.
\end{proof}

\begin{proposition}[Noncompactness]\label{modular.noncompact}
The quotient $\overline\Gamma\backslash\mathbb H$ is noncompact, despite the finite area calculated in the next section.
\end{proposition}

\begin{proof}
For $z=x+iy$, let $H(z)=\sup_{\gamma\in\Gamma}\operatorname{Im}(\gamma z)$, the maximal height of its orbit. The reduction argument shows that it is finite, and it is invariant under $\Gamma$. It is also locally bounded: if $z$ ranges over a compact subset with $\varepsilon\leq y\leq M$, the terms with $c=0$ have height at most $M$, and those with $c\ne0$ have height at most $1/\varepsilon$. This bound holds on a neighborhood of each point and descends to an open neighborhood of its quotient point. If the quotient were compact, finitely many such neighborhoods would bound $H$ globally. But for $t\geq1$, the denominator bound at $it\in\mathcal F$ gives $H(it)=t$. This is unbounded, a contradiction.
\end{proof}

The cusp neighborhood has shrinking horizontal circumference but extends to infinite hyperbolic height: for $t\geq Y$, the distance from $iY$ to $it$ is $\log(t/Y)$. Its finite area therefore does not imply compactness. This distinction is the geometric reason that analysis at the cusp will require separate estimates.
